\documentclass[11pt,a4paper]{article}
\usepackage[margin=25mm]{geometry}
\usepackage{amsmath,amssymb}
\usepackage{booktabs}
\usepackage{enumitem}
\usepackage[colorlinks=true,allcolors=black]{hyperref}
\usepackage{parskip}
\usepackage{xcolor}

\definecolor{quoteline}{gray}{0.55}

% Reviewer comment block. Replace the argument with the reviewer's verbatim text if preferred;
% the text below is our restatement of each comment.
\newenvironment{comment}
  {\begin{list}{}{\leftmargin=6mm \rightmargin=2mm}\item[]\color{quoteline}\itshape}
  {\end{list}}

\newcommand{\where}[1]{\textbf{Changes in the manuscript.} #1}

\title{\vspace{-12mm}\textbf{Response to Reviewers}\\[2mm]
\large JRESS-D-26-04700 \\
\large Similarity-Calibrated Conformal Prediction for Remaining-Useful-Life Intervals\\
\large under Operating-Regime Transfer}
\date{}

\begin{document}
\maketitle
\vspace{-10mm}

\section*{To the Editor and Reviewers}

We thank both reviewers for reports that were unusually specific and, in two places, correct about
problems we had not seen. The revision is substantial: one reviewer comment invalidated the
calibration scheme behind every number in Section~5, and we re-ran the entire experimental campaign
rather than patch around it. Five changes are worth flagging before the point-by-point replies.

\textbf{1. The calibration was not i.i.d., and every number has been recomputed.} Reviewer~1's
comment on temporal dependence was correct. Scores stacked at three monitoring points per unit
correlate at $\rho=0.908$, a design effect of 2.82, so the published calibration set of 1200 scores
carried the information of roughly 426 independent ones. The primary analysis is now unit-level,
one score per unit. The headline survives: regime-blind conformal miscovers by 0.132 and SCC
recovers coverage to within 0.011 of target, against 0.006 previously, and the change is what
theory predicts, since losing a factor of three in stacking should inflate the finite-sample floor
by $\sqrt3=1.73$, giving $0.006\times1.73=0.010$ against 0.011 observed.

\textbf{2. A second dataset, with the positive and the negative both reported.} On NASA C-MAPSS
FD002 and FD004, 509 run-to-failure engines across six operating regimes, the coverage gap falls
from 0.590 to 0.060 and from 0.567 to 0.058. The a-priori certificate, however, does not transfer:
neither physically motivated departure quantity predicts the residual gap, and the bound holds on
67 to 83\% of held-out regime pairs against 100\% on the testbed. We report this as the boundary of
the contribution rather than omitting it.

\textbf{3. The certificate claim is narrowed accordingly.} The abstract and conclusions now claim
the a-priori bound only where the governing degradation law identifies the departure quantity
$\boldsymbol\psi$, as ISO~281 does for rolling-contact fatigue and no comparable standard does for
a turbofan gas path. The dimensionless calibration, which is the part a practitioner deploys, is
claimed broadly and is now validated on external data.

\textbf{4. The actionability limit is conceded, not argued.} At $\eta=2$ the interval back-off is
1.011, so the certified lower bound is non-positive on average. The word ``graceful'' has been
removed from the manuscript and replaced by a stated operating envelope.

\textbf{5. Two errors we found ourselves during the revision, both now fixed.} The robustness table
had been computed under the old stacked calibration, which would have left two calibration schemes
side by side in one paper; it has been re-run unit-level. Its ``held-out'' bound-satisfaction rate
was in fact computed in sample; it is now a genuine 60/40 split, as the caption always claimed.

Every number in the revised manuscript is reproduced by a committed script in the public repository,
including the figures, which previously had no generator. The C-MAPSS files were verified against
the NASA release by checksum: \texttt{train\_FD002.txt} has MD5\\
\texttt{b6eaab2a6b589e5e41d43ca2f99e379b}.

\begin{center}
\small
\begin{tabular}{@{}llp{62mm}@{}}
\toprule
Comment & Outcome & Where in the revised manuscript \\
\midrule
R1.1 base-model complexity & addressed & Sec.~5.4, Table~2 \\
R1.2 deployability         & addressed & Sec.~5.8, Table~6 \\
R1.3 terminology           & accepted  & Secs.~1, 5.8, 6, 7 \\
R1.4 misspecification      & addressed & Sec.~5.7, Table~5, Sec.~3.6 \\
R1.5 temporal dependence   & conceded  & Sec.~3.4, Sec.~4.3, all of Sec.~5 \\
R2.1 second dataset        & addressed & Sec.~4.4, Sec.~5.5, Table~3, Fig.~2(b) \\
R2.2 $\boldsymbol\psi$ derivation & addressed & Appendix~B, Tables~B.1--B.2 \\
R2.3 discrete-mode baseline & addressed & Sec.~5.6, Table~4, Fig.~2(a) \\
R2.4 actionability         & conceded  & Secs.~4.3, 5.2, Fig.~3 \\
R2 minor 1 ($m$ and the norm) & addressed & Secs.~3.3, 5.3, Appendix~B \\
R2 minor 2 (bound tightness)  & conceded & Secs.~5.3, 5.4, Table~1, Fig.~4 \\
R2 minor 3 ($L$ fit)          & addressed & Sec.~3.5 \\
\bottomrule
\end{tabular}
\end{center}

\clearpage
\section*{Reviewer 1}

\subsection*{Comment 1: dependence on base-model complexity}
\begin{comment}
The results use a closed-form physics predictor. A more flexible learner might degrade the
guarantee, and the paper does not test this.
\end{comment}

\noindent\textbf{Response.} The risk is real, and the experiment shows it falls entirely on the
unscaled baseline rather than on SCC. Across four model classes, the naive coverage gap at $\eta=0$
worsens by a factor of 3.5, from 0.127 for the closed-form extrapolator to 0.449 for a random
forest, because a learner fitted flexibly to one operating regime fails harder when deployed on
another. SCC is flat at 0.007 to 0.014 across all four, and under departure the more flexible
learners are the better ones, 0.046 for a neural network against 0.102 for the closed form at
$\eta=2$, because a model fitted in the dimensionless coordinates absorbs part of the residual
departure. We also state the caveat the result carries: the drop-in property holds provided the
learned predictor is fitted in the dimensionless coordinates, since no conformal layer can repair a
point predictor that does not transfer.

We note one further finding rather than hide it. On C-MAPSS the ordering of the naive column runs
the other way, because the raw sensor levels there encode the operating regime itself, so a flexible
learner can partly infer the regime and compensate. What is invariant across both datasets is that
SCC is the better arm in every row.

\where{New second paragraph and Table~2 in Section~5.4, with the C-MAPSS contrast stated
explicitly.}

\subsection*{Comment 2: deployability and the data environment the method needs}
\begin{comment}
It is unclear how much data the diagnostic needs before it can return a usable verdict, which
matters for whether the method is deployable outside unusually rich datasets.
\end{comment}

\noindent\textbf{Response.} We have mapped it. Under exact similitude, where \emph{holds} is the
correct answer, the requirement is modest: three units per condition suffice at a unit-to-unit life
scatter of 0.10, six at 0.25 and ten at 0.50, and only at a scatter of 0.80 does the requirement
rise to about forty units. The C-MAPSS fleet, at roughly 260 units per regime, sits far inside the
feasible region.

The map also settles the status of the FEMTO result. With six bearings per condition and a
sevenfold life spread, that dataset sits in the worst corner of the map, where \emph{holds} is
returned in 0.50 to 0.00 of trials, so its \emph{indeterminate} verdict is predicted by the data
environment rather than being a quirk of the dataset.

We also report a property of our own test that the map makes visible: since the reported verdict is
the worst of three pairwise interval tests at a nominal 95\% each, the achievable rate under exact
similitude is about $0.95^3\approx0.86$, which is where the larger fleets settle. A
multiplicity-corrected verdict would raise that ceiling and is the obvious refinement; we have left
the test as submitted and stated the ceiling rather than redesign the diagnostic during revision.

\where{New opening and Table~6 in Section~5.8, with the FEMTO paragraph now referring back to the
map.}

\subsection*{Comment 3: FEMTO is described as field data}
\begin{comment}
FEMTO/PRONOSTIA is an accelerated laboratory testbed. Calling it a field benchmark overstates the
evidence.
\end{comment}

\noindent\textbf{Response.} Agreed, and corrected throughout. FEMTO is now described as an
accelerated laboratory test platform, and the word ``field'' is reserved for genuinely operational,
in-service data, which this study does not have. The heading of the corresponding results
subsection has been rewritten for the same reason. The second dataset added at Reviewer~2's request
is likewise identified as externally authored simulation rather than field data. The abstract no
longer refers to FEMTO at all, having been restructured to carry the second dataset within the
word limit; the diagnostic mechanism it exercises is described there, and the FEMTO result itself
is reported in Section~5.8.

A committed script, \texttt{scripts/r13\_terminology\_sweep.py}, audits the manuscript source for
any unqualified use of ``field'' and was run before every subsequent commit of this revision.

\where{Abstract, Section~1 scope paragraph and contribution list, Section~5.8 title and body,
Section~6 limitation~(iii), Section~7.}

\subsection*{Comment 4: what happens when the assumed physics is wrong}
\begin{comment}
The scale depends on an assumed activation energy. If that assumption is wrong, the method may fail
silently.
\end{comment}

\noindent\textbf{Response.} We have quantified both the tolerance and the detectability, and the
second half is the part that answers the ``silently''.

Tolerance separates into two cases. When only the scale is wrong and the predictor is unaffected,
the SCC gap stays at or below 0.05 out to $\lvert\Delta E/E\rvert=31\%$, remains better than not
scaling at all out to $+62\%$, and becomes worse than the unscaled baseline only at $+100\%$ (0.139
against 0.132). When a single wrong model drives both the predictor and the scale, which is the
realistic case, the method is much more fragile and is already worse than the unscaled baseline at
$-10\%$ (0.072 against 0.061).

The failure is not silent. The runtime invariance check returns \emph{holds} only at the correct
scale ($g=0.967$, interval $[0.92,1.01]$) and \emph{violated} at every misspecification tested,
including the $-10\%$ point at which the coupled case first turns against the method ($g=1.215$,
$[1.16,1.27]$). This is a capability the submitted manuscript did not claim: because the check
tests the invariance of the dimensionless life, it is sensitive to misspecification of the scale
itself as well as to a departure from similitude.

\where{New Section~5.7 with Table~5; the dual role is now stated at the end of Section~3.6.}

\subsection*{Comment 5: within-unit dependence in the calibration set}
\begin{comment}
Scores taken at several monitoring points of the same unit are not independent, so the calibration
set does not satisfy the i.i.d.\ hypothesis of the theorem.
\end{comment}

\noindent\textbf{Response.} The reviewer is right, and this was the most consequential comment in
either report. The submitted calibration stacked one score per (unit, monitoring fraction), so each
unit contributed three strongly dependent points. We measured the dependence rather than argue about
it: pairwise correlations across monitoring fractions of 0.965, 0.883 and 0.877, mean $\rho=0.908$,
design effect $1+(F-1)\rho=2.82$, so the published $n=1200$ carried an effective $n$ of about 426
against 400 units.

The primary analysis is now unit-level, one score per unit, and every number and figure in Section~5
has been recomputed on that basis. The finding is unchanged: naive conformal miscovers by 0.132
while SCC recovers to 0.011. The change in the floor, from 0.006 to 0.011, is what theory predicts
for the loss of a factor of three in stacking, $\sqrt3\times0.006=0.010$.

We also report the stronger option the comment implies. Where a guarantee is wanted simultaneously
at every monitoring point of a unit rather than at one sampled point, block conformal calibration on
the per-unit maximum score is exchangeable across units by construction; it costs conservatism, a
gap of 0.007 at $\eta=0$ and 0.173 at $\eta=2$ against 0.011 and 0.097 for single-point unit-level
calibration. We report it as the option for multi-point monitoring and state explicitly that
exchangeability is required across \emph{units}, not across time steps within a unit.

\where{New Section~3.4; the calibration scheme and the averaging convention are stated in
Section~4.3; all of Section~5, Table~1 and Figures~2 to 4 recomputed unit-level.}

\clearpage
\section*{Reviewer 2}

\subsection*{Major comment 1: no dataset the authors did not design}
\begin{comment}
The bound is never checked against an asset the authors did not design. A second, independent
dataset is needed.
\end{comment}

\noindent\textbf{Response.} We have added NASA C-MAPSS, FD002 (260 engines) and FD004 (249), six
operating regimes with every engine visiting all of them, so each regime carries roughly 260 units
against the six bearings per condition in FEMTO. The dimensionless coordinate is derived rather than
fitted, from referred (corrected) gas-turbine parameters, and it is self-checking: the sea-level
static regime returns $\theta=1.0000$ and $\delta_{\mathrm{p}}=0.9999$.

The calibration transfers. Regime-blind conformal loses a mean 0.590 of coverage across the 30
ordered regime pairs of FD002; referring the sensors to standard day is not enough on its own and
still loses 0.550; calibrating in the gas-path deviation coordinate recovers coverage to within
0.060 of the 0.90 target, a $9.8\times$ reduction, and FD004 replicates it at 0.567 to 0.058. No
target failure data is used.

The certificate does not transfer, and we report that plainly. Neither departure candidate
identifies the residual gap: the ambient referred distance correlates 0.289 with it on FD002 and
0.179 on FD004, and the healthy gas-path signature distance correlates $-0.122$ and $-0.086$; held
out, both regressions explain nothing, with $R^2$ between $-0.074$ and 0.022. The bound consequently
holds on 67 to 83\% of held-out regime pairs against 100\% on the testbed. We therefore separate the
two halves of the contribution in the abstract, in Section~5.5 and in the conclusions: the
dimensionless calibration is claimed broadly and is now validated externally; the a-priori
certificate is claimed only where the governing law identifies $\boldsymbol\psi$.

\where{New Section~4.4 (design and protocol), new Section~5.5 (results, both halves), Table~3,
Figure~2(b), and the narrowed claim in the abstract and Section~7.}

\subsection*{Major comment 2: $\boldsymbol\psi$ is never made concrete}
\begin{comment}
The departure parameter $\boldsymbol\psi$ is treated abstractly. For a real degradation law it is
not clear what it is or how a practitioner would obtain it.
\end{comment}

\noindent\textbf{Response.} A new appendix works the split through for rolling-contact fatigue,
the best-standardised degradation law in the reliability literature. Dimensional analysis on life,
load, dynamic capacity and speed gives two groups, $\Pi_1=tn$ and $\Pi_2=P/C$, and
Lundberg--Palmgren is precisely the relation between them, so the scale is
$\sigma=10^{6}(C/P)^{p}/(60n)$ in hours. ISO~281 then identifies what is left over: the modified
life $L_{nm}=a_1a_{\mathrm{ISO}}L_{10}$ has $a_{\mathrm{ISO}}=f(e_CC_u/P,\kappa)$, so for this law
the standard itself enumerates $\boldsymbol\psi$ as the lubrication regime and the contamination
group.

Evaluated on the FEMTO conditions, the appendix reproduces the $1.63\times$ characteristic-life
ratio from first principles, and it yields an instruction an engineer would not guess from the
rating-life formula: at a matched thermal state the residual departure is 0.091, while a
25\,$^\circ$C spread in operating temperature raises it to 1.060, roughly $12\times$ larger. For
bearings it is the thermal and lubrication state, not the load, that breaks similitude.

This also connects to major comment~1. Where the governing standard enumerates the residual groups,
$\boldsymbol\psi$ is identifiable and the a-priori bound applies; no equivalent enumeration exists
for a turbofan gas path, which is the most plausible reading of the C-MAPSS negative. Identifiability
of $\boldsymbol\psi$, rather than the bound itself, is the practical boundary of the method.

\where{New Appendix~B with Tables~B.1 and B.2, cross-referenced from Sections~3.2, 5.5 and 5.8.}

\subsection*{Major comment 3: comparison against discrete-mode conformal}
\begin{comment}
Treating each operating condition as a discrete mode and calibrating within it is the natural
alternative and is not compared against.
\end{comment}

\noindent\textbf{Response.} Added, with the most favourable reading given to the baseline: the
three testbed temperatures are supplied to it as known modes. The comparison is honest in both
directions.

When the target condition has its own run-to-failure history, discrete-mode calibration is
near-oracle and, once the departure is large, better than SCC (0.009 against 0.045 at $\eta=1$). We
state plainly that in that situation the discrete remedy is the right tool and the physical scale is
not needed.

When the target condition has no failure history, the discrete method has no calibration set for
that mode, falls back to the nearest mode, and lands close to regime-blind pooling, 0.078 to 0.138
against 0.093 to 0.159 pooled, while SCC holds 0.009 to 0.040, a gap smaller by $3.5\times$ to
$8.7\times$. That regime, a newly commissioned unit or a new operating point, is the one SCC exists
for.

\where{New Section~5.6 with Table~4, and a third series in Figure~2(a) showing that the nearest-mode
fallback inherits the same directional failure as regime-blind calibration.}

\subsection*{Major comment 4: are the intervals actionable at large departures}
\begin{comment}
Coverage can be preserved by widening the interval. The paper should say whether the intervals
remain useful for a maintenance decision as the departure grows.
\end{comment}

\noindent\textbf{Response.} We concede the point rather than argue it. Efficiency is now defined as
the relative back-off $b=q_T/\overline{\mathrm{RUL}}_T$, with three bands fixed in the text:
$b<0.5$ actionable, $0.5\le b<1$ degraded, $b\ge1$ degenerate, the last meaning that the certified
lower bound is non-positive on average, so the interval states only that failure has not yet
occurred.

Measured on the testbed, $b$ rises from 0.283 at $\eta=0$ to 1.011 at $\eta=2$. SCC is therefore
actionable out to $\eta=0.68$, degraded to $\eta=1.97$, and degenerate beyond, and at $\eta=2$ the
certified fraction of predicted life is $-0.011$. Coverage is still certified there, with a gap of
0.097, and we say so, but the interval carries no replacement schedule. The word ``graceful'' has
been removed from the manuscript; the operating envelope replaces it.

\where{Efficiency and the bands defined in Section~4.3; the envelope stated in Section~5.2; bands
shaded in Figure~3; ``graceful'' removed from the abstract, Section~4.3, Section~6 and Section~7.}

\subsection*{Minor comment 1: the dimension of $\boldsymbol\psi$ and the choice of norm}
\begin{comment}
Equation~(1) uses a norm of $\Delta\boldsymbol\psi$ without specifying which norm, and the tested
case appears to be one-dimensional.
\end{comment}

\noindent\textbf{Response.} Correct, and now stated. The dimension $m$ is set by the physics rather
than by the method, and the case exercised in this paper has $m=1$, a scalar weight on a single
unmodelled channel, so the norm reduces to an absolute value and its choice is immaterial. For
$m>1$ the norm must be fixed at the same time as $\boldsymbol\psi$, since $L$ is estimated in that
norm and rescales with it. Appendix~B gives the concrete two-dimensional case: carrying the
contamination group alongside the lubrication regime makes $\boldsymbol\psi$ two-dimensional, and we
say why the appendix uses the scalar $\ln\kappa$ instead.

\where{Section~3.3, the scalar instance labelled in Section~5.3, and the $m=2$ case in Appendix~B.}

\subsection*{Minor comment 2: the bound holds on 100\% of configurations, which suggests looseness}
\begin{comment}
A bound that is satisfied in every case may simply be loose. Tightness should be reported alongside
validity.
\end{comment}

\noindent\textbf{Response.} The suspicion is correct and we now report it. Across held-out
configurations the mean certified bound is 0.384 against a mean measured gap of 0.046, a margin of
0.339 with a range of 0.268 to 0.552, so the certificate is conservative by roughly a factor of
eight. We also name the two sources. The dominant one is the intercept rather than the physical
slope, since at $\delta=0$ the bound is $2a=0.278$ against a measured gap of 0.011; that intercept
is the finite-sample floor of the plug-in total-variation estimator and decays as $n^{-1/2}$ in the
number of calibration units, from 0.424 at 50 units to 0.096 at 800, a log-log slope of $-0.513$
against the $-0.500$ predicted. The second is irreducible: the swap argument of Lemma~1 bounds the
perturbation of the calibration set by twice the total-variation distance, and no amount of
calibration data removes it. The certificate is presented as a worst-case planning bound rather
than a sharp estimate.

\where{Section~5.3; a margin column in Table~1; the slack shaded and the mean margin annotated in
Figure~4.}

\subsection*{Minor comment 3: how $L$ is fitted}
\begin{comment}
The estimation of $L$ appears to rely on a small number of source pairs, which is described only in
the results.
\end{comment}

\noindent\textbf{Response.} Clarified and moved to the method, where it belongs. The design points
are configurations rather than conditions: each ordered pair of source conditions contributes one
point per distinct departure magnitude the source data exhibits, so a fleet with three conditions
supplies far more than three points. In the study reported in Section~5.3 this gives 30
(pair, departure) configurations, of which 60\% are used for the fit and the remaining 40\% are held
out to validate the resulting bound.

\where{Section~3.5.}

\section*{Reproducibility}

Every quantitative claim in the revised manuscript is reproduced by a committed script, and the
figures are now generated by one as well. The C-MAPSS files were checksum-verified against the NASA
Prognostics Data Repository release before use. We are grateful to both reviewers: the manuscript
makes a narrower claim than it did, on considerably more evidence.

\end{document}
